\documentclass[11pt]{article}
\usepackage{german,a4}
\begin{document}
\newcommand{\iu}{\int\limits_{-\infty}^{\infty}}
\newcommand{\sq}{\sqrt{2\pi}x}
\newcommand{\sqs}{\sqrt{2\pi}s}

\newcommand{\ex}{\exp\left( -i2\pi xs \right)}
\newcommand{\es}{\exp\left( i2\pi xs \right)}

\newcommand{\h}{2^{1/4}}
\newcommand{\gx}{\exp\left( - \pi x^2 \right)}
\newcommand{\gs}{\exp\left( - \pi s^2 \right)}
\newcommand{\db}{\mbox{d}}
\newtheorem{satz}{Satz}

\title{Fraktale Fouriertransformation}
\author{Pawel Wocjan}
\huge
\centerline{\bf\Huge Die Fouriertransformation}
\vfill
$$ {\cal F}[f(x)] = \iu f(x) \ex dx $$
\vfill
\LARGE
\begin{tabular}{ll}
$\bullet$ Streckung &
${\cal F}[f(x)] = \frac{1}{|a|}F\left( \frac{s}{a} \right) $ \\  & \\
$\bullet$ Linarit"at &
${\cal F}\left[ af \left( x \right) + bg \left( x \right) \right] = aF\left( s \right) + bG\left( s \right) $ \\  & \\
$\bullet$ Verschiebung &
${\cal F}[ f( x-a ) ] = e^{\left( -i2\pi as\right) F\left( s \right)} $ \\ & \\
$\bullet$ Ableitung &
${\cal F}\left[ f'\left( x \right) \right] = i2\pi s F\left( s \right) $ \\ & \\
$\bullet$ Erstes Moment &
${\cal F}[ xf(x) ] = \frac{F'(s)}{-2 \pi i} $ \\ & \\
$\bullet$ Zweites Moment &
${\cal F}[ x^2f(x) ] \frac{F''(s)}{-4 \pi} $
\end{tabular}
\vfill
\huge

\newpage
\centerline{\bf \Huge Harmonischer Oszillator}
\vfill
die zeitunabh"angige Schr"odingergleichung
$$  \left( -\frac{\hbar^2}{2m}\frac{\partial^2}{\partial x^2}+\frac{m}{2}\omega^2  x^2 \right) \varphi \left( x \right) = E\varphi\left( x \right) $$
\vfill
Substitution 
$ \xi = \frac{x}{\sigma_0} \;, \;\sigma_0 = \sqrt{\hbar / m\omega} $ liefert:
$$\frac{1}{2}\left(-\frac{\partial^2}{\partial \xi^2} + \xi^2 \right)\Phi\left( \xi \right) =  \epsilon\Phi \left(\xi\right) \: , \: \Phi\left( \xi \right)=\varphi\left(\sigma_0\xi\right) \: $$
\vfill
L"osungen nur normierbar, wenn
$$ \epsilon_n = n + \frac{1}{2} \:, \: n=0, 1,2, \ldots $$
die zugeh"origen Eigenfunktionen normiert in $\xi$ 
$$ \Phi\left(\xi\right) = \left(\sqrt{\pi}2^n n!\right)^{-1/2}H_n \left(\xi\right)e^{-\xi^2/2} $$
\vfill
die Eigenfunktionen normiert in x
$$ \varphi_n\left(x\right)=\left(\sigma_0\sqrt{\pi}2^n n!\right)^{-1/2}H_n\left(\frac{x}{\sigma_0}\right)\exp\left( -\frac{x^2}{2\sigma_0^2}\right) $$

\newpage
\centerline{\bf \Huge Der Hilbertraum $L_2\left(-\infty ,\infty\right)$ }
\vfill
$$\{ |\varphi_n\left.\right> =\varphi_n\left(x\right) \}_{n=0 \ldots \infty} $$ 
\vfill
\centerline{vollst"andiges und orthonormiertes System des }
\centerline{Hilbert\-raumes $L_2\left(-\infty ,\infty\right)$} 
\vfill
\centerline{Skalarprodukt:}
\begin{eqnarray*}
\left<f|g\right> & = & \iu \overline{f\left( x\right)} g\left( x\right) dx \: . \\
& & \\
\left<\varphi_n | \varphi_m \right> & = & \delta_{nm} \\
\sum_{n=-\infty}^{\infty}|\varphi_n\left. \right>\left<\varphi_n|\right. & = & \mbox{id} 
\end{eqnarray*}
\vfill

\newpage
\centerline{\bf \Huge Die Eigenfunktionen der FT}
\vfill
$$f''\left( x\right) + 4\pi ^2\left[ \left(2n + 1\right)/2 \pi - x^2\right] f\left( x\right) = 0$$
\vfill
\centerline{L"osungen:}
\begin{eqnarray*}
  |\Psi_n\left.\right> & = & \Psi_n \left( x \right) \\
                       & = & \frac{\h} {\sqrt{2^n n!}} H_n \left( \sq \right) \gx \\
\end{eqnarray*}
\vfill
\centerline{Fouriertransformierte der Dgl}
$$F''\left( s\right) + 4\pi ^2\left[ \left(2n + 1\right)/2 \pi - s^2\right] f\left( s\right) = 0$$
\vfill

\newpage
\vfill
\begin{eqnarray*}
{\cal F}|\Psi_n\left.\right> & = & \lambda_n \, |\Psi_n\left.\right> \\
\lambda_n & = & i^{-n}
\end{eqnarray*}
$|\Psi_n\left.\right>$ Eigenfunktionen des ${\cal F}$-Operators \\ 
Die Kernmatrix des ${\cal F}$-Operators: 
$$ \left< \Psi_m \right. \!| \, {\cal F} \, |\! \left. \Psi_n \right>_{m,n=0}^{\infty} = \mbox{Diag}\left(1,-i,-1,i,\ldots\right)$$
\vfill
\begin{eqnarray*}
  |f\left. \right> & = & \sum_{n=0}^{\infty}a_n |\Psi_n \left. \right> \\ 
  a_n & = & \left< \right. \Psi_n | f \left. \right> \\
  {\cal F}|f \left. \right> & = & \sum_{n=0}^{\infty}A_n i^{-n} |\Psi_n \left. \right> 
\end{eqnarray*}
\vfill
{\bf fraktale Fouriertransformation}
\begin{eqnarray*}
{\cal F}^a \left[{\Psi}_n \left( x \right)\right] & = & {\lambda}_n^{a} \Psi_n \left( x\right)=i^{-an}\Psi _n\left( x\right) \\
{\cal F}^{a}| f\left. \right> & = & \sum_{n=0}^{\infty}i^{-an}a_n |\Psi _n \left. \right> 
\end{eqnarray*}
\vfill

\newpage
\centerline{\bf \Huge Fraktale FT in der Optik} 
\vfill
\centerline{quadratisches GRIN-Medium} 
$$ n^2\left( \varrho \right) = \left( 1-\frac{n_2}{n_1} \varrho^2 \right) \; , \; \varrho^2 = x^2 + y^2 $$
\vfill
\centerline{Helmholtzgleichung:}
$$\nabla^2 \Psi + k^2n^2 \Psi = 0 $$ 
\vfill
\centerline{L"osungen sind die Eigenmoden:} 
$$\Psi_{n_x , n_y} \left( x,y \right) = \Psi_{n_x}\left( x \right) \Psi_{n_y}\left( y \right) $$ 
\vfill




\end{document}
 
